
il
suffit, d'après la définition de $\bar f^{*}_{U,s}$ comme adjoint à gauche du foncteur image directe, de définir un homomorphisme $f_{U,s*}(\underline{O}_X)$-linéaire
\[
  \underline{M}_{(s),U}\longrightarrow f_{U,s*}\bigl(\underline{F}|D_+(U,s)\bigr).
\]
Pour ce faire, il faut se donner pour tout ouvert $V$ de $U$ un homomorphisme :
\begin{equation}\tag{2.5.2}
  \Gamma\bigl(V,\underline{M}\bigr)_{(s)}\longrightarrow \Gamma\bigl(D_+(V,s),\underline{F}\bigr)
\end{equation}
compatible aux restrictions aux ouverts de $U$, et prendre l'homomorphisme correspondant sur les faisceaux associés. L'homomorphisme (2.5.2) est défini de la façon suivante : soit $z/s^n\in\Gamma(V,\underline{M})_{(s)}$, où, $s$ étant de degré $d$, $z\in\Gamma(V,\underline{M}_{nd})=\Gamma\bigl(D_+(V,s),\underline{F}(nd)\bigr)$; on associe alors à $z/s^n$ la section $z.\alpha_{nd}(s^n)^{-1}$ de $\underline{F}$ sur $D_+(V,s)$, en notant $\alpha_d$ l'homomorphisme canonique $\Gamma(V,\underline{B}_d)\longrightarrow \Gamma\bigl(f^{-1}(V),\underline{O}_X(d)\bigr)$ (cf. EGA II 2.6.3 et 2.6.4). On vérifie sans peine les compatibilités nécessaires, d'où par recollement l'homomorphisme (2.5.1).

Il est clair que dans le cas où $\underline{F}$ est quasi-cohérent l'homomorphisme ainsi construit est celui de EGA II (3.3.4.1).

L'homomorphisme de foncteur $\Projz\circ F_k\to \operatorname{Id}$ sur la catégorie des $\underline{O}_X$-Modules donne par passage aux catégories dérivées un homomorphisme de foncteur sur $D^+(\mathcal A)$ :
\begin{equation}\tag{2.5.3}
  \Projz\circ\mathbb R F_k\longrightarrow \operatorname{Id}_{D^+(\mathcal A)}.
\end{equation}

\textbf{Proposition 2.6.} \uline{Soient} $S$ \uline{un schéma quasi}-\uline{compact}, $\underline{E}$ \uline{un} $\underline{O}_S$\uline{-Module localement libre de type fini}, $X=\mathbf P(\underline{E})$, \uline{et} $\underline{L}^{\bullet}$ \uline{un complexe parfait sur} $X$. \uline{Alors} \uline{il existe un entier} $k_0$ \uline{tel que pour} $k\geq k_0$ \uline{les propriétés suivantes soient vérifiées} :
